\chapter{Objectives and Scope}
\label{ch:objectives}

This project aims to design, implement, and validate a reusable teaching package of practical cybersecurity labs based on EVE-NG\cite{eve_ng}. The package enables students in the ``Introduction to Cybersecurity'' course to build practical skills in ethical pentesting through realistic and automated scenarios.

\section{Specific Objectives (Measurable with KPIs)}

\subsection{O1: Design and Implement EVE-NG Labs}

\textbf{Description}: Create four functional and interconnected thematic labs.

\textbf{Measurable KPIs}:
\begin{itemize}
\item 4 fully functional .unl topologies (100\%)
\item Deployment time $\leq$~2 minutes per lab (95\% of cases)
\item VM boot success rate $\geq$~95\%
\item Fully replicable on any EVE-NG Community Edition instance with equivalent hardware
\end{itemize}

\subsection{O2: Develop Automation Scripts}

\textbf{Description}: Develop utility scripts that automate recurring lab management tasks, reducing manual effort for teachers and ensuring repeatable configurations across lab sessions.

\textbf{Measurable KPIs}:
\begin{itemize}
\item At least one functional automation script per lab covering a recurring task
\item Scripts documented and reusable across different EVE-NG installations
\item All scripts tested and verified in the project's EVE-NG environment
\end{itemize}

\subsection{O3: Develop Structured Teaching Material}

\textbf{Description}: Develop comprehensive documentation for students and teachers.

\textbf{Measurable KPIs}:
\begin{itemize}
\item 4 student lab guides (1 per lab) in PDF format --- task description,
      hints, and tool reference
\item 4 teacher resolution guides (1 per lab) in PDF format --- step-by-step
      walkthrough, expected output, and assessment rubric
\item Assessment rubric included in each resolution guide with quantifiable criteria
\item Lab setup procedure documented for teachers in each resolution guide
\end{itemize}

\subsection{O4: Validate Usability and Educational Effectiveness}

\textbf{Description}: Validate that the package meets educational and technical requirements through a pilot group representative of the target student population. The following targets are set as success criteria (see Section~\ref{sec:pilot_validation}):

\textbf{Measurable KPIs}:
\begin{itemize}
\item Average lab resolution time within the 90--120 minute target window
\item All pilot participants able to complete the lab with the provided guide
\item User satisfaction score $\geq$~4/5 (post-session survey)
\item Zero blocking technical incidents during pilot sessions
\end{itemize}

\textbf{Note}: The pilot group comprises five individuals from the GEISI student
population, potentially including participants with documented learning
difficulties (e.g.\ dyslexia), to assess the effectiveness of the
accessibility measures incorporated into the lab documentation.

\section{Definition of the Client and Final User}

This project was developed in the context of the \textit{Introduction to
Cybersecurity} course at TecnoCampus (GEISI programme). The \textbf{client}
is the course coordination body responsible for the undergraduate cybersecurity
curriculum --- concretely, the academic staff overseeing the subject --- who
identified the need for structured, reusable lab material and defined the
educational requirements that this package must satisfy.
Although TecnoCampus serves as the initial development and validation context,
the package is explicitly designed for replicability: any institution operating
an EVE-NG Community Edition instance on comparable infrastructure (e.g.,
a Proxmox-based homelab or a dedicated server) can adopt the labs, scripts,
and documentation without modification beyond local network parameters.

The \textbf{primary end users} are 4th-year GEISI students (25--30 per
edition) enrolled in the Introduction to Cybersecurity course. At that stage
of the degree, students already hold foundational knowledge of networks and
operating systems, making them ready to engage with the practical pentesting
scenarios proposed here. Their goal is to develop ethical hacking skills in a
controlled and repeatable environment.

\section{Scope Boundaries}

\textbf{Included in this bachelor's thesis:}

\begin{itemize}
    \item Design and implementation of 4 pentesting labs in EVE-NG, covering
          the four thematic domains defined by the course coordination as
          priority learning areas.
    \item Integration of real attack and defence techniques aligned with
          OWASP Top~10\cite{owasp2021}, NIST SP~800-61\cite{nist_sp80061},
          and applied cryptography and steganography, because these frameworks
          represent the industry-standard skill set expected at the 4th-year
          level.
    \item A set of utility scripts for recurring lab management tasks
          (e.g.\ bridge configuration, ISO upload, connectivity checks).
          These scripts are functional but minimal in scope: full automation
          and cross-environment generalisation were not pursued within this
          project due to validation complexity and time constraints.
          Extending this automation layer is identified as future work.
    \item Complete technical documentation for teachers and future
          maintainers, so that the package remains operational beyond the
          scope of this project.
    \item Student guides and learning objectives for each lab, directly
          supporting objective O3 and enabling autonomous lab completion.
    \item Assessment rubrics for evaluating student lab completion,
          in order to provide teachers with quantifiable grading criteria
          aligned with course competences.
\end{itemize}

The specific labs, scripts, and documentation implemented here constitute a
\emph{representative sample} of what the proposed framework supports, not an
exhaustive ceiling. The approach is intentionally extensible: additional labs
covering new attack domains, alternative network topologies, or different target
operating systems can be integrated by following the same design principles
applied throughout this project.

\textbf{Explicitly NOT included (out of scope):}

\begin{itemize}
    \item Integration with third-party Learning Management Systems (LMS) or
          Moodle, because the package is designed to be self-contained and
          deployable without dependency on any specific institutional platform.
          LMS integration would require knowledge of each institution's
          configuration, API credentials, and enrolment policies --- factors
          that vary widely between universities and are outside the technical
          scope of this project. Institutions that wish to embed the material
          in their LMS can do so independently by linking to the generated PDFs
          or importing them as resources.
    \item Development of proprietary lab management software or web
          interfaces, as EVE-NG's built-in interface is sufficient for the
          intended use and custom tooling would add maintenance overhead.
    \item Creation of automated grading systems beyond basic script
          validation, since formal assessment remains the responsibility of
          the teaching staff and falls outside the technical scope.
    \item Commercial licensing or distribution models, because the package
          is developed exclusively for academic use within the institution.
    \item Formal pedagogical evaluation studies
          (e.g., pre/post learning assessments or controlled efficacy trials),
          because such research requires longitudinal data collection across
          multiple course editions, institutional ethics approval, and a
          dedicated research methodology that falls well beyond the scope of
          an undergraduate capstone project. The pilot validation conducted in
          this thesis (Section~\ref{sec:pilot_validation}) provides indicative
          usability evidence, but it does not constitute a rigorous educational
          research study. A full pedagogical evaluation is left as future work
          for dedicated academic research.
    \item Creation of labs beyond the four specified domains, because the
          four selected topics --- network pentesting and incident response,
          web application attacks, cryptography and password analysis, and
          steganography --- cover the priority learning areas identified by
          the course coordination for the current syllabus. Extending the
          collection with additional domains (e.g., binary exploitation,
          cloud security, or IoT attacks) is entirely feasible within the
          framework but would exceed the workload constraints of a single
          bachelor's thesis. The framework's modular design explicitly
          anticipates and supports such growth in future iterations.
\end{itemize}
\newpage
\section{KPIs and Performance Indicators}

\begin{table}[H]
  \centering
  \small
  \setlength{\tabcolsep}{4pt}
  \renewcommand{\arraystretch}{1.1}
  \begin{tabularx}{\textwidth}{l l l X}
    \toprule
    \textbf{Category} & \textbf{Metric} & \textbf{Objective} & \textbf{Measurement Method} \\
    \midrule
    Technical & Lab deployment time      & $\leq$~2 min     & Automated timing \\
    Technical & VM boot success rate     & $\geq$~95\%      & System logs + validation scripts \\
    Technical & Full reset time          & $\leq$~30 s      & Scripts with timestamps \\
    Technical & Scenario reproducibility & 100\%            & Automated tests \\
    \midrule
    Educational & Lab resolution time    & 90--120 min      & Time tracking during pilot \\
    Educational & Lab completion (pilot) & All participants & Observation during pilot sessions \\
    Educational & User satisfaction      & $\geq$~4/5       & Post-use survey \\
    Educational & Key concept understanding & $\geq$~80\% correct & Assessment questionnaire \\
    \midrule
    Quality & Documentation coverage    & 100\% of features & Verification checklist \\
    Quality & Critical errors           & 0                 & Exhaustive testing \\
    Quality & EVE-NG compatibility      & 100\%             & Tests in multiple environments \\
    Quality & Interface usability       & $\geq$~4/5        & UX heuristic evaluation \\
    \bottomrule
  \end{tabularx}
  \caption{Key performance indicators for the lab teaching package}
  \label{tab:kpi_all}
\end{table}

\noindent \textbf{Note}: KPI results are drawn from the pilot validation session conducted in April--May~2026 (5~participants, Labs~1--3). Key outcomes: overall satisfaction 4.67/5, 100\% lab completion rate, mean resolution time 102~min. Full results are reported in Section~\ref{sec:pilot_validation}.

Full deliverable specifications per objective, per-category KPI breakdowns, and extended
audience analysis are provided in Appendix~F.
